\documentclass[aspectratio=169,10pt]{beamer}

\usetheme{Madrid}
\usecolortheme{default}
\setbeamertemplate{navigation symbols}{}

\usepackage[T1]{fontenc}
\usepackage[utf8]{inputenc}
\usepackage[french]{babel}
\usepackage{amsmath}
\usepackage{booktabs}
\usepackage{tabularx}
\usepackage{array}
\usepackage{ragged2e}
\usepackage{graphicx}
\usepackage{hyperref}

\title{GFlowNets et apprentissage actif pour la découverte sous budget}
\subtitle{Environnement Scrabble discret avec oracle coûteux}
\author{Yohan Zytoon, Charbel Gebrael, Yann-Olivier Jean}
\institute{IFT 3710 / IFT 6759 --- Hiver 2026}
\date{\today}

\newcolumntype{Y}{>{\RaggedRight\arraybackslash}X}
\newcommand{\code}[1]{\texttt{#1}}

\begin{document}

\begin{frame}
  \titlepage
\end{frame}

\begin{frame}{Le problème en une phrase}
\justifying
Nous voulons découvrir rapidement de \textbf{bons mots valides} dans un très
grand espace discret, alors que chaque évaluation exacte du score est
\textbf{coûteuse}.

\vspace{0.6em}
Notre banc d'essai est inspiré de Scrabble :
\begin{itemize}
  \item un candidat est une séquence de lettres,
  \item l'oracle retourne le vrai score Scrabble si le mot est valide,
  \item sinon la récompense vaut \(0\),
  \item et chaque méthode dispose d'un budget limité de requêtes oracle.
\end{itemize}
\end{frame}

\begin{frame}{Pourquoi ce problème est intéressant}
\justifying
Ce cadre est une version jouet d'un problème classique de découverte
scientifique :
\begin{itemize}
  \item l'espace de recherche est \textbf{combinatoire},
  \item la majorité des candidats sont mauvais ou invalides,
  \item l'évaluation exacte est \textbf{rare et coûteuse},
  \item on cherche donc à \textbf{dépenser intelligemment le budget}.
\end{itemize}

\vspace{0.5em}
Autrement dit, il ne suffit pas de maximiser un score. Il faut maximiser le
score \textbf{sous contrainte de coût expérimental}.
\end{frame}

\begin{frame}{Définition de la tâche}
\justifying
Chaque état est un mot de longueur maximale \(7\), représenté comme une
séquence de lettres.

\[
r(x)=
\begin{cases}
\text{ScoreScrabble}(x), & \text{si } x \text{ est un mot anglais valide},\\[0.3em]
0, & \text{sinon.}
\end{cases}
\]

\vspace{0.5em}
Dans nos expériences principales :
\begin{itemize}
  \item longueur maximale \(=7\),
  \item vérification de vocabulaire activée,
  \item budget oracle réel \(=1000\) et identique pour les méthodes comparées.
\end{itemize}
\end{frame}

\begin{frame}{Pourquoi c'est difficile}
\justifying
L'espace est énorme, mais le vrai problème est surtout la
\textbf{parcimonie de la récompense}.

\vspace{0.4em}
Même en ne gardant que les longueurs \(3\) à \(7\), on a déjà
\[
\sum_{\ell=3}^{7} 26^\ell \approx 8{,}35\times 10^9
\]
chaînes brutes possibles.

\begin{itemize}
  \item La plupart des chaînes de lettres ne sont pas des mots.
  \item Donc, la plupart des requêtes oracle retournent \(0\).
  \item Cela rend l'apprentissage plus difficile pour un modèle prédictif,
  pour un GFlowNet direct, et même pour une stratégie d'exploration naïve.
\end{itemize}

\vspace{0.5em}
Le projet porte donc sur la \textbf{recherche efficace sous récompense sparse}.
\end{frame}

\begin{frame}{Qu'est-ce que l'apprentissage actif ?}
\justifying
L'apprentissage actif est une boucle pour les cas où les labels sont coûteux.

\vspace{0.5em}
Principe :
\begin{enumerate}
  \item on part d'un petit ensemble déjà étiqueté,
  \item on entraîne un modèle substitut,
  \item on génère un pool de candidats,
  \item une fonction d'acquisition choisit les candidats les plus utiles,
  \item on interroge l'oracle uniquement sur ce sous-ensemble,
  \item on réentraîne le modèle et on recommence.
\end{enumerate}

\vspace{0.5em}
Ici, le substitut essaie de prédire quels mots valent la peine d'être évalués
pour de vrai.
\end{frame}

\begin{frame}{Qu'est-ce qu'un GFlowNet ?}
\justifying
Un \textbf{Generative Flow Network} est un modèle génératif qui apprend à
échantillonner des objets \textbf{proportionnellement à une récompense}.

\vspace{0.5em}
Intuition :
\begin{itemize}
  \item le modèle construit un objet étape par étape,
  \item il apprend une politique de génération,
  \item à convergence, les objets à forte récompense sont générés plus souvent,
  \item tout en gardant une certaine \textbf{diversité}.
\end{itemize}

\vspace{0.5em}
Dans notre contexte, c'est intéressant parce qu'on ne veut pas seulement un
seul mot optimal, mais plusieurs bons candidats.
\end{frame}

\begin{frame}{Les méthodes comparées}
\justifying
Nous comparons cinq méthodes :

\begin{enumerate}
  \item \textbf{Supervisé} :
  entraînement d'un prédicteur, puis classement d'un grand pool de candidats.
  \item \textbf{Actif} :
  boucle d'apprentissage actif standard avec substitut et acquisition.
  \item \textbf{GFlowNet} :
  GFlowNet entraîné directement sur la vraie récompense oracle.
  \item \textbf{Hybride} :
  pipeline hybride du runner de comparaison (\code{train\_hybrid.py}).
  \item \textbf{Borne supérieure} :
  scénario oracle beaucoup plus riche, non comparable en coût.
\end{enumerate}

\vspace{0.5em}
\textbf{Note :} le repo contient aussi une variante plus stricte,
\code{hybrid\_gfn\_only}, où après bootstrap l'acquisition ne voit que des
échantillons du GFlowNet. Je la présente plus loin comme variante, pas comme
la méthode du tableau principal.
\end{frame}

\begin{frame}{Vue d'ensemble des pipelines}
\justifying
Pour bien lire les résultats, il faut voir que les quatre méthodes
budget-matchées ne font pas le même travail.

\vspace{0.5em}
\begin{itemize}
  \item \textbf{Supervisé} : on entraîne une fois un prédicteur sur un dataset
  fixe, puis on classe un grand pool de candidats.
  \item \textbf{Actif} : on réentraîne un substitut à chaque ronde et on
  choisit explicitement quels candidats méritent une requête oracle.
  \item \textbf{GFlowNet} : on entraîne directement un générateur avec la vraie
  récompense oracle.
  \item \textbf{Hybride} : on entraîne un substitut, puis un GFlowNet sur ce
  substitut, puis l'acquisition choisit dans un \textbf{pool mixte}
  (échantillons GFlowNet + candidats \(n\)-gramme + mutations + recherche
  locale) quels mots envoyer à l'oracle.
\end{itemize}

\vspace{0.5em}
Le projet compare donc à la fois des \textbf{modèles prédictifs}, des
\textbf{politiques génératives}, et une \textbf{combinaison des deux}.
\end{frame}

\begin{frame}{Le substitut : ce qu'il apprend vraiment}
\justifying
Dans les résultats présentés, le substitut par défaut est un
\textbf{Gaussian Process}. Une alternative par \textbf{deep ensemble} existe
dans le code, mais ce n'est pas elle qui est utilisée ici.

\vspace{0.5em}
Le GP ne reçoit pas seulement un one-hot brut du mot. Il utilise aussi des
features de domaine :
\begin{itemize}
  \item encodage one-hot des \(7\) positions,
  \item somme des valeurs des lettres Scrabble,
  \item longueur du mot,
  \item ratio de voyelles,
  \item plausibilité \(n\)-gramme,
  \item plus longue suite de consonnes,
  \item présence ou non d'une voyelle.
\end{itemize}

\vspace{0.5em}
L'idée est simple : avec peu d'exemples positifs, ces features aident le
substitut à séparer plus vite les mots plausibles des chaînes absurdes.
\end{frame}

\begin{frame}{Apprentissage actif : config de référence}
\justifying
\small
La méthode \textbf{Actif} suit la boucle suivante :
\begin{enumerate}
  \item ensemble initial de \textbf{96} mots évalués,
  \item à chaque ronde, réentraînement du GP sur toutes les données observées,
  \item génération d'un pool principal de \textbf{8192} candidats par
  échantillonnage \(n\)-gramme,
  \item ajout de \textbf{4096} mutations construites à partir des
  \textbf{32 meilleurs} mots déjà vus,
  \item ajout d'une recherche locale guidée par le substitut
  (\textbf{512} propositions, beam width \textbf{32}, \textbf{5} étapes),
  \item filtrage des doublons, prédiction \((\mu,\sigma)\), puis acquisition.
\end{enumerate}

\vspace{0.4em}
Fonction d'acquisition utilisée ici :
\[
\mathrm{UCB}(x)=\mu(x)+\beta\,\sigma(x),
\]
avec \(\beta\) décroissant linéairement de \textbf{2,5} à \textbf{0,3}.

\vspace{0.4em}
En pratique, le code ajoute aussi un \textbf{bonus de plausibilité \(n\)-gramme}
à \(\mu(x)\) avant l'acquisition, pour éviter que l'incertitude seule ne pousse
vers des non-mots.

\vspace{0.4em}
La sélection du batch n'est pas un top-\(k\) brut : le code ajoute un bonus de
\textbf{diversité par distance de Hamming}, pour éviter de gaspiller le budget
sur des mots presque identiques.
\end{frame}

\begin{frame}{Comment fonctionne la méthode hybride}
\justifying
La méthode hybride garde la logique d'un apprentissage actif, mais y ajoute un
GFlowNet comme \textbf{générateur structuré de propositions}.

\vspace{0.5em}
Pipeline :
\begin{enumerate}
  \item collecter un petit ensemble initial de mots évalués,
  \item entraîner un modèle substitut sur ces évaluations,
  \item activer le GFlowNet seulement quand le substitut a assez de signal,
  \item entraîner le GFlowNet sur une récompense proxy venant du substitut,
  \item former un pool mixte : \textbf{GFlowNet + \(n\)-grammes + mutations +
  recherche locale},
  \item utiliser l'acquisition pour choisir lesquels envoyer à l'oracle réel.
\end{enumerate}

\vspace{0.5em}
Donc :
\begin{itemize}
  \item le \textbf{GFlowNet propose une partie du pool},
  \item l'\textbf{apprentissage actif sélectionne}.
\end{itemize}
\end{frame}

\begin{frame}{GFlowNet direct : config de référence}
\justifying
\small
Le \textbf{GFlowNet direct} utilise l'implémentation upstream du projet
\code{gflownet} avec l'objectif \textbf{trajectory balance}.

\vspace{0.5em}
Configuration principale dans notre code :
\begin{itemize}
  \item \textbf{2500} étapes d'entraînement,
  \item batch forward de \textbf{64},
  \item taux d'apprentissage \textbf{\(10^{-4}\)},
  \item \(\texttt{z\_dim}=16\), multiplicateur \(\texttt{lr\_z}=10\),
  \item \(\texttt{random\_action\_prob}=0{,}08\),
  \item politiques forward/backward en MLP de \textbf{3 couches} de
  \textbf{256} neurones.
\end{itemize}

\vspace{0.4em}
Dans le runner de comparaison multi-seed, la version comparée est légèrement
plus légère, mais le principe reste exactement le même : \textbf{vraie
récompense oracle, sans substitut}.

\vspace{0.5em}
Le point crucial est la récompense :
\begin{itemize}
  \item elle vient \textbf{directement de l'oracle réel},
  \item donc un mot invalide donne \(0\),
  \item et le modèle doit apprendre presque uniquement à partir d'un signal
  terminal très sparse.
\end{itemize}
\end{frame}

\begin{frame}{Variante : hybrid\_gfn\_only}
\justifying
\small
Le repo contient aussi une variante plus stricte :
\begin{enumerate}
  \item échantillon initial de \textbf{128} mots,
  \item phase de \textbf{bootstrap} heuristique par \(n\)-grammes pour obtenir
  assez d'exemples positifs,
  \item budget bootstrap plafonné à \textbf{20\%} du budget disponible,
  \item activation du GFlowNet dès que le bootstrap a produit assez de signal
  positif (\textbf{au moins 8} mots positifs et \textbf{3} rondes bootstrap),
  \item ensuite, à chaque ronde : GP \(\rightarrow\) GFlowNet
  \(\rightarrow\) candidats \(\rightarrow\) acquisition \(\rightarrow\) oracle.
\end{enumerate}

\vspace{0.4em}
Après le bootstrap, les candidats proposés à l'acquisition viennent du
\textbf{GFlowNet uniquement}. C'est précisément l'idée de
\code{hybrid\_gfn\_only}.

\vspace{0.4em}
\textbf{Important :} ce n'est \textbf{pas} la variante du tableau principal
ci-dessous ; c'est une ablation plus stricte présente dans le repo.
\end{frame}

\begin{frame}{Variante hybrid\_gfn\_only : récompense proxy}
\justifying
\small
Techniquement, ce GFlowNet n'est pas entraîné sur l'oracle réel mais sur une
\textbf{récompense proxy} construite à partir du substitut.

\vspace{0.4em}
Dans le code :
\[
R_{\text{base}}(x)=\exp\!\left(\beta_{\text{scale}}
\cdot \mathrm{clip}\!\left(\frac{\mu(x)}{\text{score\_max}},0,1\right)\right)
\]
avec \(\beta_{\text{scale}}=5\), puis un bonus multiplicatif de plausibilité.

\vspace{0.4em}
Autrement dit :
\begin{itemize}
  \item on utilise la \textbf{moyenne} du GP, pas UCB, pour éviter de pousser
  le GFlowNet vers des régions seulement incertaines,
  \item on transforme cette moyenne en récompense positive par une exponentielle,
  \item on amplifie les mots plausibles via le bonus \(n\)-gramme.
\end{itemize}

\vspace{0.4em}
Configuration principale :
\begin{itemize}
  \item \textbf{4000} étapes de GFlowNet par rafraîchissement,
  \item buffer replay de \textbf{1024},
  \item backward replay batch de \textbf{32},
  \item \textbf{warm start} entre deux réentraînements,
  \item réentraînement tous les \textbf{4} tours,
  \item \textbf{2048} échantillons par batch de génération, sur \textbf{3} batchs.
\end{itemize}
\end{frame}

\begin{frame}{Ce qu'on mesure}
\justifying
Nous ne regardons pas seulement le meilleur mot trouvé.

\begin{itemize}
  \item \textbf{Best} : meilleur score trouvé.
  \item \textbf{Top-10} : moyenne des dix meilleurs scores découverts.
  \item \textbf{Score moyen} : moyenne des scores des candidats évalués.
  \item \textbf{Valides (\%)} : proportion de mots valides parmi les requêtes oracle.
  \item \textbf{\(Q90\)} : nombre de requêtes nécessaires pour atteindre
  \(90\%\) du meilleur score final de la méthode.
\end{itemize}

\vspace{0.5em}
Ces métriques distinguent :
\begin{itemize}
  \item une méthode qui a un seul coup de chance,
  \item d'une méthode qui construit une bonne \textbf{frontière} de solutions.
\end{itemize}
\end{frame}

\begin{frame}{Résultats principaux}
\centering
\small
\resizebox{\linewidth}{!}{
\begin{tabular}{lcccccc}
\toprule
\textbf{Méthode} & \textbf{Best} & \textbf{Top-10} & \textbf{Score moy.} & \textbf{Valides (\%)} & \textbf{Q90} & \textbf{Fake oracle} \\
\midrule
Supervisé & 15,6 & 12,08 & 0,679 & 9,4 & 518 & 0 \\
Actif & 14,4 & 13,1 & 1,075 & 12,9 & 297 & 0 \\
GFlowNet & 12,6 & 9,36 & 0,269 & 6,3 & 377 & 0 \\
Hybride & 14,2 & 13,2 & 1,216 & 14,6 & 148 & 1\,126\,400 \\
Borne sup. & 17,4 & 15,06 & 1,605 & 80,3 & 6889 & 0 \\
\bottomrule
\end{tabular}
}

\vspace{0.8em}
\justifying
\textbf{Attention à l'interprétation :} on pourrait croire que le supervisé est
le meilleur parce qu'il a le plus grand \emph{Best}. Mais ce n'est pas le bon
critère si l'objectif est de découvrir \textbf{rapidement} un ensemble de
bons candidats valides.

\vspace{0.4em}
\textbf{Source :} ces chiffres viennent du run comparatif multi-seed complet du
repo.

\vspace{0.4em}
\textbf{Important :} la borne supérieure n'est pas budget-matchée. Elle sert
de plafond, pas de baseline équitable.
\end{frame}

\begin{frame}{Ce que le tableau ne montre pas}
\justifying
Le code suit aussi le \textbf{coût de calcul cheap}, pas seulement les requêtes
oracle réelles :
\begin{itemize}
  \item requêtes du substitut pour scorer de grands pools,
  \item requêtes du \emph{fake oracle} quand le GFlowNet s'entraîne sur le GP,
  \item et le temps total.
\end{itemize}

\vspace{0.5em}
Dans le run complet à \(5\) seeds :
\begin{itemize}
  \item \textbf{Supervisé} : \(\sim 8\,192\) requêtes cheap, \(\sim 2{,}5\) s,
  \item \textbf{Actif} : \(\sim 1{,}09\) million, \(\sim 183\) s,
  \item \textbf{Hybride} : \(\sim 1{,}93\) million, \(\sim 788\) s.
\end{itemize}

\vspace{0.5em}
Donc, \textbf{Hybride gagne en qualité de frontière}, mais pas gratuitement.
\end{frame}

\begin{frame}{Lecture des résultats}
\justifying
Le tableau raconte une histoire assez nette.

\begin{itemize}
  \item \textbf{Supervisé} obtient le meilleur \emph{Best} parmi les méthodes
  budget-matchées. Pris seul, ce chiffre pourrait faire croire que c'est la
  meilleure méthode.
  \item Mais ce \textbf{n'est pas notre objectif principal}. Nous ne cherchons
  pas seulement un coup de chance isolé : nous voulons trouver \textbf{vite}
  beaucoup de bons mots valides.
  \item \textbf{Actif} est une baseline très solide et équilibrée.
  \item \textbf{GFlowNet direct} est la méthode la plus faible sous budget.
  \item \textbf{Hybride} est la meilleure méthode si on regarde les métriques
  qui correspondent vraiment à ce but : meilleur \textbf{Top-10}, meilleur
  \textbf{score moyen}, meilleur \textbf{taux de validité}, et surtout
  meilleur \textbf{\(Q90\)}, donc la méthode la plus rapide pour atteindre des
  scores élevés.
\end{itemize}
\end{frame}

\begin{frame}{Pourquoi le GFlowNet direct souffre}
\justifying
Le GFlowNet direct doit apprendre à partir d'un signal très difficile.

\begin{enumerate}
  \item La récompense est très sparse : beaucoup de zéros.
  \item Le budget oracle est limité, donc peu de signal utile.
  \item Le modèle doit découvrir à la fois la validité des mots et leur score.
\end{enumerate}

\vspace{0.5em}
Ce n'est donc pas que l'idée des GFlowNets est mauvaise. C'est plutôt que
\textbf{dans ce régime de budget serré}, le problème est trop dur pour un
GFlowNet entraîné uniquement sur l'oracle réel.
\end{frame}

\begin{frame}{Pourquoi l'hybride marche mieux}
\justifying
La méthode hybride donne au GFlowNet un rôle plus réaliste.

\begin{itemize}
  \item Le substitut fournit un signal plus dense et moins brutal que l'oracle.
  \item Le GFlowNet sert à générer des candidats divers et prometteurs.
  \item L'acquisition dépense les requêtes oracle sur les meilleurs choix.
\end{itemize}

\vspace{0.5em}
Cela explique bien les chiffres :
\begin{itemize}
  \item meilleur \textbf{Top-10},
  \item meilleur \textbf{score moyen},
  \item meilleur \textbf{taux de validité},
  \item et surtout meilleur \textbf{\(Q90\)} parmi les méthodes budget-matchées,
  \item mais avec un \textbf{coût de calcul cheap} beaucoup plus élevé.
\end{itemize}
\end{frame}

\begin{frame}{Message principal du projet}
\justifying
\begin{block}{Conclusion}
Sous budget oracle serré, le \textbf{GFlowNet direct} ne suffit pas.
Le meilleur usage du GFlowNet dans notre projet est de le combiner avec une
boucle d'apprentissage actif, où il sert de générateur de candidats.
\end{block}

\vspace{0.5em}
En d'autres termes :
\begin{itemize}
  \item si on veut tester l'idée pure, le GFlowNet seul est intéressant mais
  faible dans ce cadre,
  \item si on veut être efficace expérimentalement, l'hybride est plus fort,
  \item mais ce gain doit être lu avec son \textbf{coût de calcul}.
\end{itemize}
\end{frame}

\begin{frame}{Conclusion}
\justifying
Les leçons les plus importantes sont :

\begin{itemize}
  \item la \textbf{parcimonie de la récompense} est le principal obstacle,
  \item l'\textbf{apprentissage actif} est très fort pour allouer un budget rare,
  \item le \textbf{GFlowNet} est plus utile comme générateur de propositions
  que comme solveur direct sous budget,
  \item dans le tableau principal du repo, la meilleure \textbf{qualité de
  frontière} vient de la \textbf{méthode hybride},
  \item mais cette amélioration a un \textbf{coût computationnel} nettement plus
  élevé.
\end{itemize}
\end{frame}

\end{document}
